% !TeX root = ../教材题库_全部条目_转写工作稿.tex
% Source: XJ-B1-C06-HW-P231-Y0446
% 第6章 / 6.3 统计图表 / 教材 p.231 / PDF 第238页
% OCR未校对：公式与文字均可能有误；未完成数学TeX转写；禁止直接发给学生
% 原题及相邻原解截图（核对用）：../题目原图/XJ-B1-C06-HW-P231-Y0446-01.webp
% 本文件只在转写完成后作为题干使用；OCR 条目尚不是可靠题干。
\begin{bankitem}{习题6.3 第1题}{第6章 / 6.3 统计图表 / 教材 p.231 / PDF 第238页}{XJ-B1-C06-HW-P231-Y0446}{OCR未校对：公式与文字均可能有误；未完成数学TeX转写；禁止直接发给学生}
\begin{ocrdraft}
[教材p.231 / PDF第238页]
\par
1.
\par
小敏一星期的总开支分布如图（1)所示，一星期的食品开支如图（2）所示，
\par
则小敏一星期的鸡蛋开支占总开支的百分比为多少？
\par
其他5％
\par
开支/元
\par
120
\par
食品开支
\par
100
\par
100
\par
储蓄30％
\par
30％
\par
80
\par
60
\par
40
\par
40
\par
通信开支5％
\par
20
\par
日常开支20％
\par
娱乐开支10％
\par
0
\par
鸡蛋
\par
牛奶
\par
肉类
\par
蔬菜
\par
其他食品种类
\par
(1)
\par
（2）
\par
（第1题）
\end{ocrdraft}
\end{bankitem}
